\chapter{Data and Knowledge Base}
\label{ch:du-lieu}

An advisory assistant is only as trustworthy as the data behind it. This chapter describes the book data source, the cleaning pipeline, the datasets it produces, and the knowledge base used for retrieval. Every number in this chapter is taken from the manifest and quality report generated by the processing pipeline itself \src{15}, \src{16}.

\section{Data Source}

The raw data is the \emph{Tiki Books Dataset}, publicly shared on Kaggle \src{13}, version 4, under the CC0-1.0 licence (which permits free use). The team downloaded the dataset on 24 August 2026. It consists of two tables:

\begin{itemize}
    \item a book information file: 1,796 rows, each a book title with its name, author, publisher, price, rating, page count and category;
    \item a review file: 141,281 rows, each a buyer review with a star rating.
\end{itemize}

To ensure that every run uses exactly the same copy of the data, the team recorded the SHA-256 hash of the original archive. A hash is a string computed from a file's content; changing even a single character of the file changes the hash. Before each processing run, the pipeline recomputes the hash and compares it with the recorded value; if they do not match, it stops immediately.

It must be stressed that this is a \textbf{historical snapshot}, not live data from the marketplace. Prices and ratings in this report are the values at the time the dataset was collected, not current prices.

\section{Data Preparation Pipeline}

The raw data cannot be used as is: it contains duplicate rows, empty reviews, reviews pointing to non-existent books, and private information written in by buyers themselves. The team built an eight-step pipeline, shown in Figure~\ref{fig:xu-ly-du-lieu}. The entire pipeline runs automatically and produces the same result on every run, thanks to a fixed random seed in the sampling step.

\hinh{xu-ly-du-lieu}{The eight data preparation steps, from the raw files to the database}{fig:xu-ly-du-lieu}

\textbf{Steps 1 and 2: download and verify.} The file is downloaded under a temporary name and is renamed to its official name only after the hash matches. As a result, a partially downloaded file can never be used by mistake. The pipeline also checks that the archive contains exactly the two tables with exactly the expected columns.

\textbf{Step 3: cleaning and de-duplication.} This step is applied to the entire raw dataset, before sampling. Table~\ref{tab:lam-sach} shows the number of records removed and the reasons.

\begin{table}[!htbp]
\caption{Cleaning results on the entire raw dataset}
\label{tab:lam-sach}
\cobang
\begin{tabularx}{\linewidth}{YR{3.2cm}}
\toprule
\textbf{Item} & \textbf{Records}\\
\midrule
Book titles in the raw file & 1,796\\
\quad Removed as duplicate identifiers & 19\\
\textbf{Book titles after cleaning} & \textbf{1,777}\\
\midrule
Reviews in the raw file & 141,281\\
\quad Removed for empty content & 33,451\\
\quad Removed for not linking to any valid book title & 15,659\\
\quad Removed as duplicate identifiers & 15,416\\
\textbf{Reviews after cleaning} & \textbf{76,755}\\
\bottomrule
\end{tabularx}
\end{table}

Thus nearly half of the original reviews (45.7\%) are unusable. Skipping this step would affect both the statistics on reviews and the assistant's answers.

\textbf{Step 4: masking private information.} Buyers sometimes write phone numbers, email addresses or links into their reviews. The pipeline replaces such information with generic tags. Across the whole dataset, 118 phone numbers, 76 links and 21 email addresses were masked. Buyer identifiers are also removed entirely and never enter the processed data.

\textbf{Step 5: profile-based sampling.} From the clean data, the pipeline produces datasets of different sizes for different purposes (Section~\ref{sec:ho-so}). Books are selected in rotation across categories so that the sample spans many categories. Reviews for each book are selected per star level to retain both praise and criticism.

\textbf{Step 6: writing the output.} Each dataset consists of exactly four files: a book file, a review file, a manifest (recording the origin, version, counts and hash of each file) and a quality report (recording the number of records in, out, removed and masked, and the fill rate of each field).

\textbf{Step 7: quality gate.} Before use, a dataset must pass a series of checks. The pipeline stops if any of the following occurs: duplicate identifiers; reviews not linked to any book; leftover phone numbers, emails or links; records missing a required identifier; records not from the correct source marketplace; counts or hashes not matching the manifest; any of the four files missing.

\textbf{Step 8: loading into the database.} Data is loaded in a single transaction: either everything is written or nothing is. After loading, the pipeline reads the data back and compares it with the manifest once more. Reloading the same dataset creates no duplicate records; loading a different dataset into a database that already holds an older one is rejected.

\section{Generated Datasets}
\label{sec:ho-so}

From the same clean data, the pipeline produces three datasets, as shown in Table~\ref{tab:ho-so}.

\begin{table}[!htbp]
\caption{The three datasets generated from the clean data}
\label{tab:ho-so}
\cobang
\begin{tabularx}{\linewidth}{P{2.6cm}R{1.8cm}R{2.0cm}Y}
\toprule
\textbf{Dataset} & \textbf{Books} & \textbf{Reviews} & \textbf{Purpose}\\
\midrule
Test set & 24 & 115 & Automated testing; small so it runs quickly. At most 5 reviews per book.\\
Evaluation set & 200 & 1,773 & Used for all experiments in this report and for the default deployment. At most 10 reviews per book.\\
Full set & 1,777 & 76,755 & All clean data, unsampled. Not stored in the code repository due to its size; can be regenerated when needed.\\
\bottomrule
\end{tabularx}
\end{table}

All experimental results in this report use the \textbf{evaluation set}. The rest of this section describes that dataset.

\subsection{Information for Each Book Title}

Each book title has the fields listed in Table~\ref{tab:truong-sach}. The last column gives the proportion of books with data in that field; these figures come from the quality report.

\begin{table}[!htbp]
\caption{Information fields of a book title in the evaluation set}
\label{tab:truong-sach}
\cobang
\begin{tabularx}{\linewidth}{P{3.6cm}YR{2.6cm}}
\toprule
\textbf{Field} & \textbf{Meaning} & \textbf{Fill rate}\\
\midrule
Title & Full name as on the sales page, including edition notes & 100\%\\
Sale price, list price & In VND, at the time of collection & 100\%\\
Rating & Average score on a 5-star scale & 100\%\\
Category & Book group according to the sales site's classification & 100\%\\
Popularity & Popularity index provided by the source & 98.0\%\\
Page count & Number of pages listed on the sales page & 87.0\%\\
Author & List of authors & 86.5\%\\
Publisher & Publisher name & 84.5\%\\
Seller & Not present in the raw data & 0\%\\
\bottomrule
\end{tabularx}
\end{table}

Two points deserve attention. First, the \emph{popularity} field in the raw data carries a name resembling ``quantity sold'', but there is no way to verify that it reflects actual sales. The system therefore treats it only as a relative popularity index and never calls it sales. Second, the data contains no seller information and no review timestamps at all, so the system cannot answer questions about sellers or about trends over time.

\subsection{Some Descriptive Statistics}

Table~\ref{tab:mo-ta} summarises the evaluation set. The figures were computed by the team directly from the two data files.

\begin{table}[!htbp]
\caption{Descriptive statistics of the evaluation set}
\label{tab:mo-ta}
\cobang
\begin{tabularx}{\linewidth}{YR{5.4cm}}
\toprule
\textbf{Quantity} & \textbf{Value}\\
\midrule
Book titles & 200\\
Categories & 74\\
Lowest -- highest sale price & 9,600 -- 590,000 VND\\
Middle sale price (median) & about 109,600 VND\\
Mean rating of the book titles & 4.72 out of 5\\
Page count (over 174 books with data), median & 296 pages\\
\midrule
Reviews & 1,773\\
Reviews per book & from 1 to 10\\
Review length, median & 65 characters\\
5-star reviews & 827 (46.6\%)\\
4-star reviews & 379 (21.4\%)\\
3-star reviews & 224 (12.6\%)\\
2-star reviews & 124 (7.0\%)\\
1-star reviews & 219 (12.4\%)\\
\bottomrule
\end{tabularx}
\end{table}

Most reviews are 4 or 5 stars, consistent with the general pattern on e-commerce sites. Nevertheless, thanks to sampling reviews per star level, the dataset still retains 343 one- and two-star reviews (19.3\%). These are an important source of information for questions such as ``what do buyers complain about''. Reviews are generally short: half of them are no longer than 65 characters.

Among the 200 records, one is not a book (a pack of milk misfiled under the book category in the raw data). The team did not delete this record, in order to keep exactly what the source provides. Instead, the system has a runtime filter that excludes this record from search results and statistics.

\section{Knowledge Base}
\label{sec:kho-tri-thuc}

\subsection{Why a Knowledge Base Is Needed}

The catalogue and reviews can answer questions about price, rating and purchasing experience. But readers also ask about the work itself: what the book is about, who the author is, what the original title is, who translated it. The descriptions in the catalogue are usually marketing copy, not a reliable source for such questions. Letting the model answer from memory offers no guarantee of correctness.

The team therefore built a small knowledge base: for each work, the team found an authoritative source (the publisher's page, the author's page, or a library catalogue), then \emph{wrote in its own words} a short note stating the verifiable information. The knowledge base contains no book content, does not copy publishers' marketing text, contains no buyer reviews, and contains none of the questions or answers of the evaluation set.

\subsection{Coverage}

Each book title in the evaluation set is placed in exactly one of three groups, as shown in Figure~\ref{fig:do-phu}.

\hinh{do-phu-tri-thuc}{Extent to which the 200 book titles can be matched against the knowledge base}{fig:do-phu}

\begin{itemize}
    \item \textbf{20 books with a source confirming the work.} An authoritative source confirms which work it is and by whom. The knowledge base has exactly 20 notes for these 20 works; the full list is in Appendix~\ref{app:kho-tri-thuc}.
    \item \textbf{17 records are bundles of multiple works.} These are book sets or combos of several volumes sold together. Assigning a whole bundle to a single work would be wrong, so the team created no notes for them.
    \item \textbf{163 books without a matched source yet.} In the first round of building the knowledge base, the team had not yet located and checked sources for these books.
\end{itemize}

The knowledge base therefore covers only 10\% of the book titles. This is a major limitation of the project and directly affects the results: for the remaining 180 books, every question about the content of the work must receive an ``insufficient evidence'' response.

\subsection{Distinguishing Works from Editions}

A subtle but important detail is the distinction between a \emph{work} and an \emph{edition}. A work is the content written by the author. An edition is a specific printing, with its own page count, cover, year of printing and translator. Authoritative sources usually describe an English edition, whereas the books in the catalogue are Vietnamese editions.

The team adopted the convention that knowledge base notes state only work-level information (author, subject, original title), not edition-level information (page count, year of printing). Table~\ref{tab:vi-du-an-ban} gives three examples of discrepancies between source and catalogue; without separating these two levels, the system would produce contradictory figures.

\begin{table}[!htbp]
\caption{Examples of discrepancies between authoritative sources and the catalogue}
\label{tab:vi-du-an-ban}
\cobang
\begin{tabularx}{\linewidth}{P{4.2cm}YY}
\toprule
\textbf{Work} & \textbf{Authoritative source states} & \textbf{Catalogue states}\\
\midrule
Days at the Morisaki Bookshop & National Library of Vietnam catalogue: 177 pages & 180 pages\\
Steve Jobs (biography) & English-language publisher: 656 pages & 758 pages (Vietnamese edition)\\
Start-up Nation & Two authors: Dan Senor and Saul Singer & Only one author listed\\
\bottomrule
\end{tabularx}
\end{table}

In the third example, the catalogue omits one author. The knowledge base note keeps both authors, as in the authoritative source. When a user asks about the author, the system answers according to the source and states clearly where the information comes from.

\subsection{Storage and Retrieval}

Table~\ref{tab:kho-tri-thuc} summarises the parameters of the knowledge base.

\begin{table}[!htbp]
\caption{Parameters of the knowledge base used in the experiments}
\label{tab:kho-tri-thuc}
\cobang
\begin{tabularx}{\linewidth}{P{5.4cm}Y}
\toprule
\textbf{Parameter} & \textbf{Value}\\
\midrule
Number of sources & 20\\
Number of indexed passages & 20 (one passage per source)\\
Source types & Publishers, author pages, national library, public book catalogues\\
Text embedding model & text-embedding-3-small, 1,536-dimensional vectors\\
Search method & Hybrid search: keyword and semantic\\
Storage & Same PostgreSQL database as the catalogue\\
Source verification date & 9 September 2026\\
\bottomrule
\end{tabularx}
\end{table}

Each note is stored with a version, the source access date and a content hash. When a turn begins, the system records the knowledge base version in use and keeps that version fixed until the answer is complete. As a result, every citation in an answer always points to exactly the content the system read, even if the knowledge base is updated afterwards.

Read access to the knowledge base is also controlled: only properly authenticated users with data-read permission can perform retrieval.

\section{Chapter Summary}

A public dataset with a clear licence passes an eight-step checked pipeline to give an evaluation set of 200 titles and 1,773 reviews. The knowledge base has 20 notes from authoritative sources, covering 10\% of the titles. Its main limits, missing seller and time data and a small knowledge base, are handled by stating clearly when evidence is insufficient.
